\documentclass[11pt]{article}

\title{Stone-Weierstrass for Functions Vanishing at Infinity}
\author{UlamLib seed note}
\date{}

\newtheorem{definition}{Definition}[section]
\newtheorem{lemma}[definition]{Lemma}
\newtheorem{theorem}[definition]{Theorem}
\newtheorem{corollary}[definition]{Corollary}

\begin{document}
\maketitle

\section{Scope}

This note packages a non-compact Stone-Weierstrass target for
$C_0(X, \mathbb{R})$ and related function algebras. The intended Lean
destination is \texttt{ULAM/Analysis/StoneWeierstrassC0.lean}.

\section{Core definitions}

\begin{definition}
Let $X$ be a locally compact Hausdorff space. The algebra $C_0(X, \mathbb{R})$
consists of continuous real-valued functions on $X$ that vanish at infinity.
\end{definition}

\begin{definition}
A subalgebra $A \subseteq C_0(X, \mathbb{R})$ separates points if for every
distinct pair $x,y \in X$ there exists $f \in A$ such that $f(x) \neq f(y)$.
\end{definition}

\section{Seed theorem cluster}

\begin{lemma}
If a real subalgebra of $C_0(X, \mathbb{R})$ is closed under absolute value,
then it is closed under pointwise supremum and infimum.
\end{lemma}

\begin{theorem}[Real $C_0$ Stone-Weierstrass]
Let $X$ be locally compact Hausdorff and let
$A \subseteq C_0(X, \mathbb{R})$ be a subalgebra that separates points and
vanishes nowhere. Then $A$ is dense in $C_0(X, \mathbb{R})$ for the supremum
norm.
\end{theorem}

\begin{corollary}
Compactly supported continuous functions can be uniformly approximated by
elements of any dense subalgebra satisfying the hypotheses above.
\end{corollary}

\section{Formalization backlog}

\begin{enumerate}
\item Define the separation API for subalgebras of $C_0(X, k)$.
\item Add lattice and closure lemmas in the real case.
\item Prove the real Stone-Weierstrass theorem for $C_0$.
\item Extend to complex and star-subalgebra variants.
\item Package common approximation corollaries.
\end{enumerate}

\end{document}
